\documentclass[11pt,oneside]{article}

\usepackage[T1]{fontenc}
\usepackage[utf8]{inputenc}
\usepackage{lmodern}
\usepackage{microtype}
\usepackage{geometry}
\geometry{margin=1in,headheight=15pt}
\usepackage{amsmath,amssymb,amsthm,mathtools}
\usepackage{booktabs,longtable,tabularx,array}
\usepackage{enumitem}
\usepackage{xcolor}
\usepackage{xurl}
\usepackage{hyperref}
\usepackage{listings}
\usepackage{fancyhdr}
\usepackage{titlesec}
\usepackage{needspace}

\definecolor{deepblue}{RGB}{24,74,128}
\definecolor{deepgreen}{RGB}{16,105,73}
\definecolor{deepviolet}{RGB}{90,55,135}
\definecolor{deepred}{RGB}{150,45,45}
\definecolor{softgray}{RGB}{247,248,250}
\definecolor{rulegray}{RGB}{95,104,116}

\hypersetup{
  colorlinks=true,
  linkcolor=deepblue,
  citecolor=deepgreen,
  urlcolor=deepblue,
  pdftitle={Alaoglu--Erdos Beyond Report A: S-unit Secant Heights and Reciprocal Cyclotomic Barriers},
  pdfauthor={OpenAI GPT-5.6 Pro},
  pdfsubject={New conditional consequences and no-go theorems for the Alaoglu--Erdos conjecture}
}
\urlstyle{same}

\pagestyle{fancy}
\fancyhf{}
\lhead{Alaoglu--Erd\H{o}s beyond Report A}
\rhead{Secant heights and reciprocal filters}
\cfoot{\thepage}
\renewcommand{\headrulewidth}{0.4pt}

\titleformat{\section}{\Large\bfseries}{\thesection}{0.75em}{}
\titleformat{\subsection}{\large\bfseries}{\thesubsection}{0.75em}{}
\titleformat{\subsubsection}{\normalsize\bfseries}{\thesubsubsection}{0.75em}{}

\setlist[itemize]{leftmargin=2em,itemsep=0.3em,topsep=0.45em}
\setlist[enumerate]{leftmargin=2.2em,itemsep=0.3em,topsep=0.45em}
\setlist[description]{leftmargin=3.2cm,labelsep=0.6cm,itemsep=0.35em,topsep=0.45em}
\setlength{\parindent}{0pt}
\setlength{\parskip}{0.55em}
\setlength{\emergencystretch}{3em}

\newtheorem{theorem}{Theorem}[section]
\newtheorem{proposition}[theorem]{Proposition}
\newtheorem{lemma}[theorem]{Lemma}
\newtheorem{corollary}[theorem]{Corollary}
\newtheorem{conjecture}[theorem]{Conjecture}
\newtheorem{criterion}[theorem]{Criterion}
\theoremstyle{definition}
\newtheorem{definition}[theorem]{Definition}
\newtheorem{experiment}[theorem]{Computation}
\newtheorem{target}[theorem]{Research target}
\theoremstyle{remark}
\newtheorem{remark}[theorem]{Remark}
\newtheorem{warning}[theorem]{Caution}

\newcommand{\N}{\mathbb N}
\newcommand{\Nzero}{\mathbb N_0}
\newcommand{\Z}{\mathbb Z}
\newcommand{\Q}{\mathbb Q}
\newcommand{\R}{\mathbb R}
\newcommand{\Gm}{\mathbb G_m}
\newcommand{\vp}[2]{v_{#1}(#2)}
\newcommand{\rad}{\operatorname{rad}}
\newcommand{\ord}{\operatorname{ord}}
\newcommand{\supp}{\operatorname{supp}}
\newcommand{\lcm}{\operatorname{lcm}}
\newcommand{\height}{h}
\newcommand{\eps}{\varepsilon}
\newcommand{\proved}{\textcolor{deepgreen}{\textsf{[proved here]}}}
\newcommand{\classical}{\textcolor{deepblue}{\textsf{[classical input]}}}
\newcommand{\computed}{\textcolor{deepviolet}{\textsf{[computation]}}}
\newcommand{\openstatus}{\textcolor{deepred}{\textsf{[open target]}}}

\newenvironment{resultbox}
 {\begin{quote}\hrule\smallskip\noindent}
 {\par\smallskip\hrule\end{quote}}

\lstdefinestyle{py}{
  language=Python,
  basicstyle=\ttfamily\scriptsize,
  backgroundcolor=\color{softgray},
  frame=single,
  rulecolor=\color{rulegray},
  breaklines=true,
  columns=fullflexible,
  keepspaces=true,
  showstringspaces=false,
  tabsize=4
}

\title{\bfseries Alaoglu--Erd\H{o}s Beyond Report A\\[0.35em]
\Large $S$-unit secant heights, primitive residual cancellation,\\
fixed-order extrapolation, and reciprocal cyclotomic barriers}
\author{OpenAI GPT-5.6 Pro}
\date{22 August 2026}

\begin{document}
\maketitle

\begin{abstract}
The Alaoglu--Erd\H{o}s integer-exponent conjecture asks whether
\[
  2^x,3^x\in\Z \quad\Longrightarrow\quad x\in\Z
\]
for real $x$.  The attached Report A already develops the exact conditional monoid,
finite verification, determinant hierarchies, $p$-adic assignment methods, rational and
algebraic output spectra, interpolation, operator, Mahler, moment, and canonical-product
approaches.  This continuation deliberately does not repeat those lines.

Assuming a least nonintegral solution $\beta$, with $M=2^\beta$ and $A=3^\beta$, we
normalize continued-fraction convergents $r_j/s_j\to\beta$ into rational $S$-unit points
\[
 (u_j,v_j)=\left(\frac{2^{r_j}}{M^{s_j}},\frac{3^{r_j}}{A^{s_j}}\right)
           =(2^{\eps_j},3^{\eps_j})\longrightarrow(1,1),
 \qquad \eps_j=r_j-s_j\beta.
\]
Corvaja--Zannier height theory then gives an exact linear asymptotic for the reduced height
of the rational secant $(u_j-1)/(v_j-1)$.  The coefficient is an explicit valuation
functional $\kappa(M,A;\beta)>0$.  Consequently the reduced numerator and denominator of the
secant are each $\exp((\kappa+o(1))s_j)$, while its error as an approximation to
$\log_3 2$ is only polynomial in $s_j$.  This produces a sharp ``height exponent zero''
barrier.  It also yields an exact asymptotic for the total cross-residual gcd and shows that
its entire exponential part is supported on the fixed primes dividing $6MA$; all auxiliary
prime cancellation is subexponential.

The same mechanism rules out every fixed-order Richardson/Prouhet extrapolation of toric
secants: order $m$ improves the analytic error to $\Theta(|\eps_j|^m)$, but the reduced
height remains $\Theta(s_j)$.  A separate algebraic argument proves that two distinct monic
reciprocal polynomials of the same degree can agree at infinity to at most half their degree.
For equal-totient cyclotomic pairs this is an exact universal ceiling.  A scan through index
$50000$ finds $665095$ equal-totient pairs, confirms the ceiling, and identifies only $24$
sharp pairs in that range.

These results do not prove the conjecture.  They do, however, eliminate a broad new class of
fixed-complexity secant, finite-difference, common-residual, and reciprocal-cyclotomic
amplifiers, and isolate the surviving requirement: a growing-complexity construction whose
\emph{primitive} height is saved by output-specific cancellation at the structural primes.
\end{abstract}

\tableofcontents
\clearpage

\section{Executive answer and claim ledger}
\label{sec:executive}

\subsection{Outcome}

No complete proof is claimed here.  The four-exponentials obstruction remains real: a
nonintegral solution would make the four algebraic numbers $2,3,2^\beta,3^\beta$ occupy the
forbidden $2\times2$ configuration predicted by the Four Exponentials Conjecture.  What is
new is a conditional transformation of such a solution into a near-identity rank-two
$S$-unit orbit on which the Subspace Theorem has unusually sharp consequences.

The main conclusions are the following.

\begin{enumerate}
\item \textbf{Exact secant height.}  For convergents $r_j/s_j$ to $\beta$, the rational
secants
\[
 \sigma_j=\frac{u_j-1}{v_j-1}
 =\frac{A^{s_j}(2^{r_j}-M^{s_j})}
        {M^{s_j}(3^{r_j}-A^{s_j})}
\]
satisfy
\[
 h(\sigma_j)=\kappa(M,A;\beta)s_j+o(s_j),
\]
where $\kappa$ is the explicit positive valuation functional in
Theorem~\ref{thm:kappa}.  \proved

\item \textbf{Primitive numerator and denominator.}  If $\sigma_j=P_j/Q_j$ in lowest
terms, then
\[
 \log P_j=\log Q_j=\kappa s_j+o(s_j).
\]
At the same time
\[
 \sigma_j-\log_3 2=\Theta(\eps_j),
 \qquad s_j^{-C}\ll|\eps_j|<s_j^{-1}.
\]
Thus the approximation exponent relative to reduced height tends to zero.  \proved

\item \textbf{Exact cross-residual cancellation.}  With
\[
 X_j=A^{s_j}|2^{r_j}-M^{s_j}|,
 \qquad
 Y_j=M^{s_j}|3^{r_j}-A^{s_j}|,
 \qquad
 G_j=\gcd(X_j,Y_j),
\]
one has
\[
 \log G_j=(\log M+\log A-\kappa)s_j+o(s_j).
\]
After removing the fixed primes dividing $6MA$, the remaining common divisor is
$\exp(o(s_j))$.  Hence every linear-rate cancellation in $G_j$ is structural, not supplied
by fresh primes.  \proved

\item \textbf{Fixed-order extrapolation barrier.}  Put
\[
 \phi(t)=\frac{2^t-1}{3^t-1},\qquad \phi(0)=\log_3 2,
\]
and, for fixed $m\ge1$,
\[
 R_m(t)=\sum_{i=1}^m(-1)^{i-1}\binom mi\phi(it).
\]
Then $R_m(\eps_j)\in\Q$,
\[
 R_m(\eps_j)-\phi(0)=\Theta(\eps_j^m),
 \qquad h(R_m(\eps_j))=\Theta(s_j),
\]
and again the height-normalized approximation exponent is zero.  This covers every fixed
finite-difference order, not merely the first secant.  \proved

\item \textbf{Reciprocal half-contact ceiling.}  If $P,Q\in\Z[Z]$ are distinct monic
reciprocal polynomials of common degree $D$ and
\[
 \frac{P(Z)}{Q(Z)}=1+cZ^{-K}+O(Z^{-K-1}),\qquad c\ne0,
\]
then $K\le\lfloor D/2\rfloor$.  Equality forces $P-Q$ to be a single middle monomial.
For $P=\Phi_n,Q=\Phi_m$ with $\varphi(n)=\varphi(m)$ this bounds the first unequal
M\"obius layer by $\varphi(n)/2$.  \proved

\item \textbf{Finite cyclotomic reconnaissance.}  Among all $665095$ pairs
$3\le n<m\le50000$ with $\varphi(n)=\varphi(m)$, the maximum contact ratio is $1/2$,
attained by $24$ pairs.  Every sharp pair found belongs to one of two infinite elementary
families or is the exceptional pair $(3,4)$.  No global classification is claimed.
\computed
\end{enumerate}

\begin{resultbox}
\textbf{Central message.}
The near-identity orbit supplies arbitrarily accurate analytic contact, but fixed rational
observables retain linear primitive height.  The Subspace Theorem does not merely fail to
finish the proof: it identifies exactly where the lost size goes.  Any successful continuation
must use a family whose complexity grows and whose reduced height enjoys cancellation not
available to any fixed rational function on $\Gm^2$.
\end{resultbox}

\subsection{Evidence labels}

\begin{description}
\item[\proved] A complete conventional proof is included in this report.  No Lean kernel
claim is made.
\item[\classical] A named theorem from the literature is used, chiefly the
Corvaja--Zannier $S$-unit height theorem and a standard Baker lower bound for a linear form in
two logarithms.
\item[\computed] A finite Python computation is reproduced in Appendix~\ref{app:code}; it
is not a formal proof of an infinite statement.
\item[\openstatus] A precise sufficient target or an unresolved extension.
\end{description}

\subsection{Public-source reconnaissance}

The model-comparison background was checked only to calibrate the competitive context, not
to import mathematical assertions.  Anthropic's public launch page for Claude Fable 5 is
\url{https://www.anthropic.com/news/claude-fable-5-mythos-5}.  Its 28 July 2026 research
post, \emph{Discovering cryptographic weaknesses with Claude}, states that Fable 5 had
resolved the Jacobian Conjecture earlier in July; see
\url{https://www.anthropic.com/research/discovering-cryptographic-weaknesses}.  I found no
public primary document disclosing the claimed Alaoglu--Erd\H{o}s idea.  The mathematics below
is therefore independent of any undisclosed rival construction.

\subsection{Nonduplication boundary}

Report A already treats the exact solution monoid, candidate counting, finite scans,
Four/Six Exponentials reductions, prime-log-product independence, arbitrary-node and mixed
semigroup determinants, $p$-adic assignment cascades, Loewner and matrix-coefficient
certificates, rational contact rigidity, geometric Newton interpolation, Mahler and natural
boundary diagnostics, operator formulations, moments, Pad\'e defects, carry languages,
canonical products, and several denominator barriers.  This continuation uses only the
minimal conditional normalization from that corpus.  Its new object is the convergent-normalized
$S$-unit orbit and the primitive height of rational observables on that orbit.

\section{Conditional setup and the near-identity orbit}
\label{sec:setup}

\subsection{The least generator}

Assume the conjecture is false.  Let $\beta>0$ be its least nonintegral solution and put
\begin{equation}
 M=2^\beta\in\N,\qquad A=3^\beta\in\N.
 \label{eq:MA}
\end{equation}
The inherited exact classification gives all solutions as $n+k\beta$ with
$n,k\in\Nzero$.  Only three elementary consequences are needed here:

\begin{enumerate}
\item $\beta$ is irrational, indeed transcendental;
\item $M$ is not a pure power of $2$, and $A$ is not a pure power of $3$;
\item $\log M=\beta\log2$ and $\log A=\beta\log3$.
\end{enumerate}

Let $r_j/s_j$ be the ordinary continued-fraction convergents to $\beta$, with $s_j>0$, and
set
\begin{equation}
 \eps_j=r_j-s_j\beta.
 \label{eq:epsilon}
\end{equation}
Then $\eps_j\ne0$, $\eps_j\to0$, and
\begin{equation}
 |\eps_j|<\frac1{s_j}.
 \label{eq:cf-upper}
\end{equation}
A lower bound for the nonzero linear form
\[
 r_j\log2-s_j\log M=\eps_j\log2
\]
in two logarithms of algebraic numbers gives constants $C_0,C_1>0$, depending only on $M$,
such that
\begin{equation}
 |\eps_j|\ge C_0s_j^{-C_1}
 \label{eq:baker-lower}
\end{equation}
for all $j$.  This is the only quantitative transcendence estimate used below.

\subsection{The normalized \texorpdfstring{$S$}{S}-unit point}

Define
\begin{equation}
 u_j=\frac{2^{r_j}}{M^{s_j}},\qquad
 v_j=\frac{3^{r_j}}{A^{s_j}}.
 \label{eq:uv}
\end{equation}
By~\eqref{eq:MA} and~\eqref{eq:epsilon}, as positive real numbers,
\begin{equation}
 u_j=2^{\eps_j},\qquad v_j=3^{\eps_j}.
 \label{eq:real-path}
\end{equation}
Thus $(u_j,v_j)\to(1,1)$ along the analytic one-parameter subgroup
$t\mapsto(2^t,3^t)$, while each coordinate is a rational number supported on the fixed set of
primes dividing $6MA$.

\begin{lemma}[Multiplicative independence]
\label{lem:uv-independent}
For every $j$, the rational numbers $u_j$ and $v_j$ are multiplicatively independent.
\end{lemma}

\begin{proof}
Suppose $u_j^a v_j^b=1$ with $a,b\in\Z$.  Taking the real logarithm and using
$\eps_j\ne0$ gives
\[
 a\log2+b\log3=0.
\]
Hence $2^a3^b=1$, so unique factorization gives $a=b=0$.
\end{proof}

The points lie in the finitely generated subgroup
\begin{equation}
 \Gamma=\langle(2,3),(M^{-1},A^{-1})\rangle\subset\Gm^2(\Q).
 \label{eq:Gamma}
\end{equation}
This is precisely the setting of Corvaja--Zannier's height theorem.

\subsection{Coordinate heights}

Write
\begin{equation}
 M=2^mM_0,\quad 2\nmid M_0,
 \qquad
 A=3^aA_0,\quad 3\nmid A_0.
 \label{eq:basefree}
\end{equation}
The nonintegrality of $\beta$ implies $M_0>1$ and $A_0>1$.  Since $r_j/s_j\to\beta$,
we have $r_j>ms_j$ and $r_j>as_j$ for all sufficiently large $j$.  The fractions
\[
 u_j=\frac{2^{r_j-ms_j}}{M_0^{s_j}},\qquad
 v_j=\frac{3^{r_j-as_j}}{A_0^{s_j}}
\]
are then in lowest terms.  Consequently
\begin{align}
 h(u_j)&=s_j\log M_0+O(|\eps_j|),
 \label{eq:hu}\\
 h(v_j)&=s_j\log A_0+O(|\eps_j|).
 \label{eq:hv}
\end{align}
In particular both heights tend linearly to infinity even though the real values tend to $1$.
This archimedean/nonarchimedean split is the source of the method.

\section{The exact projective-height slope}
\label{sec:kappa}

For a rational projective point, use the absolute logarithmic Weil height
\[
 h(1:u:v)=\sum_{w}\log\max\{1,|u|_w,|v|_w\},
\]
with the usual normalized absolute values on $\Q$.

\begin{theorem}[Exact projective-height slope]
\label{thm:kappa}
Define
\begin{equation}
 \boxed{
 \kappa(M,A;\beta)=\log\operatorname{lcm}(M_0,A_0).}
 \label{eq:kappa}
\end{equation}
For all sufficiently large $j$ one has the exact formula
\begin{equation}
 h(1:u_j:v_j)=s_j\log\operatorname{lcm}(M_0,A_0)
 +\max\{0,\eps_j\log2,\eps_j\log3\}.
 \label{eq:projective-asymptotic}
\end{equation}
In particular,
\begin{equation}
 \kappa=\log\operatorname{lcm}(M_0,A_0)
 \ge\max\{\log M_0,\log A_0\}>0.
 \label{eq:kappa-positive}
\end{equation}
\proved
\end{theorem}

\begin{proof}
Put $m_p=v_p(M)$ and $a_p=v_p(A)$.  At a finite prime $p$,
\begin{align*}
 v_p(u_j)&=r_j\mathbf1_{p=2}-s_jm_p,\\
 v_p(v_j)&=r_j\mathbf1_{p=3}-s_ja_p.
\end{align*}
For all sufficiently large $j$, the finite local contribution is therefore
$s_j\max\{v_p(M_0),v_p(A_0)\}\log p$.  Summing over $p$ gives
$s_j\log\operatorname{lcm}(M_0,A_0)$.  At the archimedean place,
$u_j=2^{\eps_j}$ and $v_j=3^{\eps_j}$, giving precisely the maximum in
\eqref{eq:projective-asymptotic}.  Finally, a nonintegral candidate cannot have
$M_0=1$ or $A_0=1$, so the lcm is greater than one.
\end{proof}

\begin{remark}[What $\kappa$ measures]
The real magnitudes of $u_j$ and $v_j$ contain essentially no height: both approach one.
The coefficient $\kappa$ is the surviving denominator slope after all common projective
content among $1,u_j,v_j$ has been removed.  It is therefore already a \emph{primitive}
quantity.  Raw numerator sizes are larger and can be misleading.
\end{remark}

\begin{remark}[A useful defect identity]
Set
\begin{equation}
 \lambda=\log M+\log A-\kappa.
 \label{eq:lambda}
\end{equation}
Then the local expression simplifies to
\begin{equation}
 \lambda=m\log2+a\log3+\log\gcd(M_0,A_0)\ge0.
 \label{eq:lambda-local}
\end{equation}
Thus the structural-prime contributions are explicit, while the remaining term is the
logarithm of the common prime part of $M_0$ and $A_0$.  The slope may be zero.
\end{remark}

\section{\texorpdfstring{$S$}{S}-unit secants and exact primitive height}
\label{sec:secant}

\subsection{Corvaja--Zannier input}

We use the following classical consequence of the Subspace Theorem.

\begin{theorem}[Corvaja--Zannier, 2005]
\label{thm:CZ}
Let $\Gamma\subset\Gm^2(\Q)$ be finitely generated.  As multiplicatively independent
$(u,v)\in\Gamma$ satisfy $\max\{h(u),h(v)\}\to\infty$,
\begin{equation}
 h\!\left(\frac{u-1}{v-1}\right)\sim h(1:u:v).
 \label{eq:CZ-secant}
\end{equation}
More generally, for a fixed nonconstant $F\in\Q(X,Y)$, outside a finite union of translates
of proper subtori, $h(F(u,v))$ is bounded below by a positive constant times
$\max\{h(u),h(v)\}$.
\classical
\end{theorem}

The cited paper states~\eqref{eq:CZ-secant} as Corollary 1 and the rational-function statement
as its Main Theorem.  The exceptional translates are equations $X^cY^d=w$.

\begin{lemma}[Coset avoidance]
\label{lem:coset-avoidance}
The sequence $(u_j,v_j)$ meets every fixed proper translate of a subtorus of $\Gm^2$ only
finitely many times.
\end{lemma}

\begin{proof}
A one-dimensional translate has an equation $X^cY^d=w$ with $(c,d)\ne(0,0)$.  If infinitely
many $(u_j,v_j)$ lie on it, then
\[
 (2^c3^d)^{\eps_j}=w
\]
for infinitely many $j$.  Letting $j\to\infty$ gives $w=1$.  Since $2^c3^d\ne1$ and
$\eps_j\ne0$, the left side is never $1$, a contradiction.
\end{proof}

\subsection{The principal secant theorem}

Define
\begin{equation}
 \sigma_j=\frac{u_j-1}{v_j-1}.
 \label{eq:sigma}
\end{equation}
The signs of numerator and denominator agree, so $\sigma_j>0$.

\begin{theorem}[Exact primitive secant law]
\label{thm:secant-law}
Under the counterexample hypothesis,
\begin{equation}
 h(\sigma_j)=\kappa s_j+o(s_j).
 \label{eq:sigma-height}
\end{equation}
If $\sigma_j=P_j/Q_j$ in lowest terms with $P_j,Q_j\in\N$, then
\begin{equation}
 \log P_j=\log Q_j=\kappa s_j+o(s_j).
 \label{eq:PQ-height}
\end{equation}
Finally,
\begin{equation}
 \sigma_j=\log_3 2
 +\tfrac12\log_3 2\,(\log2-\log3)\,\eps_j
 +O(\eps_j^2).
 \label{eq:sigma-expansion}
\end{equation}
\proved
\end{theorem}

\begin{proof}
Lemma~\ref{lem:uv-independent}, equations~\eqref{eq:hu}--\eqref{eq:hv}, and
Theorem~\ref{thm:CZ} give
\[
 h(\sigma_j)\sim h(1:u_j:v_j).
\]
Theorem~\ref{thm:kappa} proves~\eqref{eq:sigma-height}.  Since
$\sigma_j\to\log_3 2\in(0,\infty)$, the reduced numerator and denominator differ by only
a bounded multiplicative factor.  Their maximum has logarithm $h(\sigma_j)$, proving
\eqref{eq:PQ-height}.

For~\eqref{eq:sigma-expansion}, write $a=\log2$, $b=\log3$ and use
\[
 \sigma_j=\frac{e^{a\eps_j}-1}{e^{b\eps_j}-1}.
\]
Division of the two Taylor series gives
\[
 \frac{e^{at}-1}{e^{bt}-1}
 =\frac ab+\frac{a(a-b)}{2b}t+O(t^2).
\]
\end{proof}

The limit $\log_3 2$ is transcendental: if it were algebraic irrational,
Gelfond--Schneider applied to $3^{\log_3 2}=2$ would give a contradiction; if it were
rational, unique factorization would give a nontrivial equality between powers of $2$ and
$3$.

\begin{corollary}[Height exponent zero]
\label{cor:height-zero}
There are constants $c,C>0$ such that
\begin{equation}
 s_j^{-C}\ll
 \left|\sigma_j-\log_3 2\right|
 \ll s_j^{-1},
 \label{eq:sigma-poly}
\end{equation}
and
\begin{equation}
 \frac{-\log\left|\sigma_j-\log_3 2\right|}
      {h(\sigma_j)}\longrightarrow0.
 \label{eq:height-exponent-zero}
\end{equation}
\proved
\end{corollary}

\begin{proof}
The coefficient of $\eps_j$ in~\eqref{eq:sigma-expansion} is nonzero.  Apply
\eqref{eq:cf-upper}, \eqref{eq:baker-lower}, and~\eqref{eq:sigma-height}.
\end{proof}

\begin{resultbox}
The secants do not merely have ``large denominators.''  Their \emph{reduced} numerator and
denominator have an exact exponential rate $\kappa$.  Continued fractions improve the real
error only polynomially.  Thus no irrationality-measure argument can be won by taking these
secants at face value: the approximation exponent measured in their true height is zero.
\end{resultbox}

\subsection{Directional secants}

The principal secant is one member of a toric family.  Let
$\mathbf c=(c_1,c_2)$ and $\mathbf d=(d_1,d_2)$ be two linearly independent vectors in
$\Z^2$, and put
\[
 U_j=u_j^{c_1}v_j^{c_2},\qquad
 V_j=u_j^{d_1}v_j^{d_2}.
\]
Then
\[
 U_j=(2^{c_1}3^{c_2})^{\eps_j},\qquad
 V_j=(2^{d_1}3^{d_2})^{\eps_j},
\]
and $U_j,V_j$ are multiplicatively independent.  Corvaja--Zannier therefore gives
\begin{equation}
 h\!\left(\frac{U_j-1}{V_j-1}\right)
 \sim h(1:U_j:V_j)=\kappa_{\mathbf c,\mathbf d}s_j+o(s_j),
 \label{eq:directional-height}
\end{equation}
where
\begin{equation}
 \kappa_{\mathbf c,\mathbf d}=
 \sum_p\max\{0,L_{\mathbf c}(p),L_{\mathbf d}(p)\}\log p,
 \label{eq:kappa-cd}
\end{equation}
with
\[
 L_{\mathbf c}(p)=c_1m_p+c_2a_p-
 \beta(c_1\mathbf1_{p=2}+c_2\mathbf1_{p=3})
\]
and similarly for $\mathbf d$.  The secant tends to
\[
 \log_{2^{d_1}3^{d_2}}(2^{c_1}3^{c_2})
\]
whenever the denominator is nonzero, again with only polynomial error and linear primitive
height.  Thus changing toric directions does not evade the barrier.

\section{Cross-residual gcds: exact total rate and auxiliary-prime extinction}
\label{sec:gcd}

\subsection{Raw residuals}

Put
\begin{equation}
 E_{2,j}=2^{r_j}-M^{s_j},\qquad
 E_{3,j}=3^{r_j}-A^{s_j}.
 \label{eq:residuals}
\end{equation}
Both are nonzero and have the sign of $\eps_j$.  Define the positive cross-multiplied
integers
\begin{equation}
 X_j=A^{s_j}|E_{2,j}|,
 \qquad
 Y_j=M^{s_j}|E_{3,j}|.
 \label{eq:XY}
\end{equation}
Then $\sigma_j=X_j/Y_j$.

\begin{lemma}[Raw size]
\label{lem:raw-size}
One has
\begin{equation}
 \log X_j=\log Y_j=(\log M+\log A)s_j+O(\log s_j).
 \label{eq:raw-size}
\end{equation}
\end{lemma}

\begin{proof}
Since $2^{r_j}=M^{s_j}e^{\eps_j\log2}$,
\[
 |E_{2,j}|=M^{s_j}|e^{\eps_j\log2}-1|.
\]
Equations~\eqref{eq:cf-upper} and~\eqref{eq:baker-lower} imply
$\log|e^{\eps_j\log2}-1|=O(\log s_j)$.  Multiplication by $A^{s_j}$ gives the first formula;
the second is identical.
\end{proof}

\subsection{The exact total gcd law}

\begin{theorem}[Primitive cancellation identity]
\label{thm:total-gcd}
Let
\begin{equation}
 G_j=\gcd(X_j,Y_j).
 \label{eq:G}
\end{equation}
Then
\begin{equation}
 \boxed{
 \log G_j=(\log M+\log A-\kappa)s_j+o(s_j).}
 \label{eq:G-asymptotic}
\end{equation}
\proved
\end{theorem}

\begin{proof}
Writing $\sigma_j=P_j/Q_j$ in lowest terms gives
$P_j=X_j/G_j$ and $Q_j=Y_j/G_j$.  Subtract~\eqref{eq:PQ-height} from
\eqref{eq:raw-size}.
\end{proof}

This theorem quantifies a phenomenon that raw formulas obscure.  The numerator and
denominator in~\eqref{eq:sigma} have the explicit nonnegative common-divisor rate $\lambda$,
which may vanish.  Any positive part of that rate is not free arithmetic information: after
division, the surviving height still has the positive linear slope $\kappa$.

\subsection{No exponential common divisor from new primes}

Let $S$ be the finite set of rational primes dividing $6MA$.  Define the prime-to-$S$ common
residual divisor
\begin{equation}
 g_j^{(S)}=
 \prod_{p\notin S}p^{\min\{v_p(E_{2,j}),v_p(E_{3,j})\}}.
 \label{eq:gS}
\end{equation}

\begin{theorem}[Auxiliary-prime extinction]
\label{thm:aux-gcd}
For every $\delta>0$ and all sufficiently large $j$,
\begin{equation}
 \log g_j^{(S)}<\delta s_j.
 \label{eq:gS-small}
\end{equation}
Equivalently,
\begin{equation}
 \log g_j^{(S)}=o(s_j).
 \label{eq:gS-o}
\end{equation}
\proved
\end{theorem}

\begin{proof}
At $p\notin S$, the denominators $M^{s_j}$ and $A^{s_j}$ are $p$-adic units, so
\[
 v_p(u_j-1)=v_p(E_{2,j}),\qquad
 v_p(v_j-1)=v_p(E_{3,j}).
\]
The generalized-gcd form of Corvaja--Zannier's Corollary 1 says that for multiplicatively
independent $(u_j,v_j)$ the total common local proximity of $u_j-1$ and $v_j-1$ is less than
$\delta\max\{h(u_j),h(v_j)\}$ outside a finite set.  Restricting the nonnegative gcd sum to
$p\notin S$ can only decrease it.  Equations~\eqref{eq:hu}--\eqref{eq:hv} then give
\eqref{eq:gS-small}; since $\delta$ is arbitrary,~\eqref{eq:gS-o} follows.
\end{proof}

\begin{corollary}[Where the exponential gcd lives]
\label{cor:gcd-support}
The prime-to-$S$ part of $G_j$ is $\exp(o(s_j))$.  Hence the full linear rate
$\log M+\log A-\kappa$ in~\eqref{eq:G-asymptotic} is carried by the fixed structural primes
in $S$.
\end{corollary}

\begin{proof}
Outside $S$, the factors $M^{s_j}$ and $A^{s_j}$ in~\eqref{eq:XY} contribute no valuation,
so the prime-to-$S$ part of $G_j$ is exactly $g_j^{(S)}$.
\end{proof}

\begin{remark}[Finite families]
Take finitely many nonproportional toric directions
$\mathbf c_1,\ldots,\mathbf c_t\in\Z^2$ and the residuals
$u_j^{c_{i,1}}v_j^{c_{i,2}}-1$.  The prime-to-$S$ gcd of all cleared residuals is bounded by
the gcd of any multiplicatively independent pair, and is therefore $\exp(o(s_j))$.  Thus no
fixed toric jet can acquire an exponential auxiliary common divisor.
\end{remark}

\begin{resultbox}
The total cross-residual gcd is large, but its linear part is completely structural.  This
rules out a common proposed escape hatch: manufacture many near-identity residuals and hope
that fresh primes clear the denominator.  Fresh-prime clearing is subexponential on every
fixed finite toric configuration.
\end{resultbox}

\section{Fixed-order Richardson extrapolation}
\label{sec:richardson}

The first secant has error $O(\eps_j)$.  It is natural to cancel successive Taylor terms by
combining secants at $\eps_j,2\eps_j,\ldots,m\eps_j$.  The next theorem shows exactly why every
fixed order still fails.

\subsection{The analytic coefficients}

Set
\begin{equation}
 \phi(t)=\frac{2^t-1}{3^t-1}
 \quad(t\ne0),
 \qquad
 \phi(0)=\rho:=\log_3 2.
 \label{eq:phi}
\end{equation}
This is real analytic at $0$.  Write
\begin{equation}
 \phi(t)=\sum_{k\ge0}c_kt^k.
 \label{eq:phi-series}
\end{equation}

\begin{lemma}[No vanished Taylor coefficient]
\label{lem:ck-nonzero}
Every coefficient $c_k$ is nonzero.
\end{lemma}

\begin{proof}
Put $a=\log2$ and $b=\log3$.  From
\[
 \frac{e^{at}-1}{e^{bt}-1}
 =\frac{e^{at}-1}{bt}\cdot\frac{bt}{e^{bt}-1}
\]
and the Bernoulli expansion $z/(e^z-1)=\sum_{n\ge0}B_nz^n/n!$, the coefficient $c_k$ is a
homogeneous rational polynomial in $a,b$ divided by $b$.  Its term of highest degree in $a$
is
\[
 \frac{a^{k+1}}{b(k+1)!},
\]
so the polynomial is not identically zero.  After substituting $b=\theta a$ with
$\theta=\log_2 3$ transcendental, a nonzero rational polynomial in $\theta$ cannot vanish.
\end{proof}

\subsection{The extrapolants}

For fixed $m\ge1$, define
\begin{equation}
 R_m(t)=\sum_{i=1}^m(-1)^{i-1}\binom mi\phi(it).
 \label{eq:Rm}
\end{equation}
At $t=\eps_j$ every summand is rational:
\begin{equation}
 \phi(i\eps_j)=\frac{u_j^i-1}{v_j^i-1}\in\Q.
 \label{eq:phi-rational}
\end{equation}
Thus $R_m(\eps_j)\in\Q$.

\begin{theorem}[Exact fixed-order barrier]
\label{thm:richardson}
For every fixed $m\ge1$ there are nonzero constants $C_m$ and positive constants
$c_m',C_m'$ such that, as $j\to\infty$,
\begin{align}
 R_m(\eps_j)-\rho
 &=C_m\eps_j^m+O(\eps_j^{m+1}),
 \label{eq:Rm-error}\\
 c_m's_j\le h(R_m(\eps_j))&\le C_m's_j,
 \label{eq:Rm-height}\\
 \frac{-\log|R_m(\eps_j)-\rho|}{h(R_m(\eps_j))}&\longrightarrow0.
 \label{eq:Rm-zero}
\end{align}
\proved
\end{theorem}

\begin{proof}
The binomial moment identity gives
\begin{align*}
 \sum_{i=1}^m(-1)^{i-1}\binom mi&=1,\\
 \sum_{i=1}^m(-1)^{i-1}\binom mi i^k&=0
 \quad(1\le k<m),\\
 \sum_{i=1}^m(-1)^{i-1}\binom mi i^m&=(-1)^{m-1}m!.
\end{align*}
Substitution into~\eqref{eq:phi-series}, together with
Lemma~\ref{lem:ck-nonzero}, proves~\eqref{eq:Rm-error} with
$C_m=(-1)^{m-1}m!c_m\ne0$.

Now define the fixed rational function
\begin{equation}
 F_m(X,Y)=\sum_{i=1}^m(-1)^{i-1}\binom mi\frac{X^i-1}{Y^i-1}
 \in\Q(X,Y).
 \label{eq:Fm}
\end{equation}
It is nonconstant: along $(X,Y)=(2^t,3^t)$ it tends to $\rho$ as $t\to0$ but tends to $0$
as $t\to+\infty$.  Corvaja--Zannier's Main Theorem gives a positive linear lower bound for
$h(F_m(u_j,v_j))$ outside finitely many subtorus translates.  Lemma~\ref{lem:coset-avoidance}
removes those exceptions.  The standard height upper bound for a fixed rational map gives
the other inequality in~\eqref{eq:Rm-height}.  Finally,~\eqref{eq:baker-lower},
\eqref{eq:cf-upper}, and~\eqref{eq:Rm-error} show that the numerator in~\eqref{eq:Rm-zero} is
$O(\log s_j)$, while the denominator is $\Theta(s_j)$.
\end{proof}

\begin{remark}[What increasing the fixed order accomplishes]
The construction genuinely improves contact: order $m$ produces $m$th-order error.  The
failure is arithmetic, not analytic.  Each fixed $F_m$ still has linear primitive height on
the $S$-unit orbit.  Consequently no finite $m$, however large, changes the height exponent.
\end{remark}

\begin{warning}[Growing order is the surviving loophole]
Theorem~\ref{thm:richardson} is uniform only after $m$ is fixed.  If $m=m_j\to\infty$, the
rational function, its exceptional subtori, its coefficients, and its height constants all
change.  That is not a technical omission: growing complexity is precisely the only regime
not eliminated by the theorem.  Any proof along this line must control the \emph{reduced}
height of the growing family, not merely a common denominator before cancellation.
\end{warning}

\section{Reciprocal filters and the half-contact theorem}
\label{sec:reciprocal}

A second natural amplifier uses products of cyclotomic polynomials or other reciprocal
polynomials to cancel many leading powers at infinity.  Reciprocity imposes a rigid ceiling
that appears not to have been used in Report A.

\subsection{The algebraic ceiling}

A polynomial $P\in\Z[Z]$ of degree $D$ is \emph{reciprocal} if
\begin{equation}
 Z^DP(Z^{-1})=P(Z).
 \label{eq:reciprocal}
\end{equation}

\begin{theorem}[Reciprocal half-contact]
\label{thm:half-contact}
Let $P,Q\in\Z[Z]$ be distinct monic reciprocal polynomials of the same degree $D$ and with
constant term $1$.  Suppose
\begin{equation}
 \frac{P(Z)}{Q(Z)}=1+cZ^{-K}+O(Z^{-K-1}),
 \qquad c\ne0.
 \label{eq:contact-infinity}
\end{equation}
Then
\begin{equation}
 2K\le D.
 \label{eq:half-bound}
\end{equation}
Equality holds if and only if $D$ is even and
\begin{equation}
 P(Z)-Q(Z)=cZ^{D/2}.
 \label{eq:middle-monomial}
\end{equation}
Equivalently,~\eqref{eq:middle-monomial} characterizes equality.
\proved
\end{theorem}

\begin{proof}
Let $H=P-Q$.  Because $P$ and $Q$ are monic with constant term $1$, the leading and constant
coefficients of $H$ vanish.  Reciprocity gives
\[
 Z^DH(Z^{-1})=H(Z),
\]
so the coefficient $h_i$ of $Z^i$ equals $h_{D-i}$.  If $d$ is the largest index with
$h_d\ne0$, the smallest such index is $D-d$.  Hence $D-d\le d$, or $d\ge D/2$.
Since $Q(Z)\sim Z^D$, the contact order in~\eqref{eq:contact-infinity} is $K=D-d$, proving
$2K\le D$.  Equality forces the smallest and largest nonzero indices of $H$ both to be
$D/2$, so $H$ is a single middle monomial.  The converse is immediate.
\end{proof}

\subsection{A two-scale integer determinant}

For fixed $P,Q$ as above, define
\begin{equation}
 \Delta_s(P,Q)=P(N^s)Q(M^s)-Q(N^s)P(M^s)\in\Z,
 \qquad N:=A=M^\theta,
 \label{eq:two-scale-det}
\end{equation}
where $\theta=\log_2 3$ and $A=M^\theta$ under the counterexample hypothesis.

\begin{proposition}[Exact growth of a fixed reciprocal filter]
\label{prop:reciprocal-growth}
With $K$ as in~\eqref{eq:contact-infinity},
\begin{equation}
 \log|\Delta_s(P,Q)|
 =\bigl(D(1+\theta)-K\bigr)s\log M+O(1).
 \label{eq:det-growth}
\end{equation}
In particular, by Theorem~\ref{thm:half-contact},
\begin{equation}
 \log|\Delta_s(P,Q)|
 \ge D\left(\theta+\frac12\right)s\log M+O(1).
 \label{eq:det-growth-lower}
\end{equation}
Thus every fixed reciprocal filter produces a rapidly growing nonzero integer rather than a
subunit contradiction.  \proved
\end{proposition}

\begin{proof}
Write $R=P/Q$.  Then
\[
 \Delta_s=Q(N^s)Q(M^s)\bigl(R(N^s)-R(M^s)\bigr).
\]
Since $Q(Z)=Z^D(1+O(Z^{-1}))$ and
$R(Z)=1+cZ^{-K}+O(Z^{-K-1})$,
\begin{align*}
 Q(N^s)Q(M^s)&=M^{D(1+\theta)s}(1+o(1)),\\
 R(N^s)-R(M^s)&=-cM^{-Ks}(1+o(1)).
\end{align*}
Multiplication gives~\eqref{eq:det-growth};~\eqref{eq:det-growth-lower} follows from
$K\le D/2$.
\end{proof}

\begin{remark}
The proposition addresses a fixed filter evaluated deeper and deeper on the conditional
orbit.  It does not rule out a sequence $P_s,Q_s$ of increasing degree with output-dependent
coefficient cancellation.  Once again, the only surviving regime is growing complexity.
\end{remark}

\section{Cyclotomic specialization}
\label{sec:cyclotomic}

For $n>1$, the cyclotomic polynomial $\Phi_n$ is monic, reciprocal, and has constant term
$1$.  If $\varphi(n)=\varphi(m)=D$, Theorem~\ref{thm:half-contact} applies to
$P=\Phi_n,Q=\Phi_m$.

\subsection{The M\"obius-layer description}

The standard factorization gives, as a Laurent expansion at infinity,
\begin{equation}
 \Phi_n(Z)=Z^{\varphi(n)}
 \prod_{d\mid n}(1-Z^{-d})^{\mu(n/d)}.
 \label{eq:cyclotomic-infinity}
\end{equation}
For equal totients define
\begin{equation}
 c_d(n,m)=
 \mathbf1_{d\mid n}\mu(n/d)-
 \mathbf1_{d\mid m}\mu(m/d).
 \label{eq:cd}
\end{equation}
Then
\begin{equation}
 \frac{\Phi_n(Z)}{\Phi_m(Z)}
 =\prod_{d\ge1}(1-Z^{-d})^{c_d(n,m)}.
 \label{eq:cyclotomic-ratio}
\end{equation}
If
\begin{equation}
 K(n,m)=\min\{d\ge1:c_d(n,m)\ne0\},
 \label{eq:Knm}
\end{equation}
then the first unequal M\"obius layer is exactly the contact order at infinity.

\begin{corollary}[Cyclotomic half-contact]
\label{cor:cyclotomic-half}
For distinct $n,m>1$ with $\varphi(n)=\varphi(m)=D$,
\begin{equation}
 K(n,m)\le\left\lfloor\frac D2\right\rfloor.
 \label{eq:cyclotomic-half}
\end{equation}
Equality forces
$\Phi_n-\Phi_m$ to be a single middle monomial.  \proved
\end{corollary}

\subsection{Sharp infinite families}

Two elementary families attain the ceiling.

For $a\ge2$, let $h=2^{a-2}$.  Then
\begin{align}
 \Phi_{2^a}(Z)&=Z^{2h}+1,\\
 \Phi_{3\cdot2^{a-1}}(Z)&=Z^{2h}-Z^h+1,
\end{align}
so their difference is $Z^h$ and $K=D/2=h$.

For $a\ge1$, let $h=3^{a-1}$.  Then
\begin{align}
 \Phi_{3^a}(Z)&=Z^{2h}+Z^h+1,\\
 \Phi_{2\cdot3^a}(Z)&=Z^{2h}-Z^h+1,
\end{align}
so their difference is $2Z^h$ and again $K=D/2=h$.  The pair $(3,4)$ is another sharp
example:
\[
 \Phi_3(Z)-\Phi_4(Z)=Z.
\]

\subsection{Finite scan through \texorpdfstring{$50000$}{50000}}

\begin{experiment}[Equal-totient cyclotomic scan]
\label{exp:scan}
The script in Appendix~\ref{app:code} enumerates all pairs
\[
 3\le n<m\le50000,\qquad \varphi(n)=\varphi(m),
\]
computes $K(n,m)$ from~\eqref{eq:cd}, and checks the ratio $K/\varphi(n)$.
It produced:
\begin{center}
\begin{tabular}{lr}
\toprule
Quantity & Value\\
\midrule
Equal-totient pairs tested & $665095$\\
Maximum $K/\varphi(n)$ & $1/2$\\
Pairs attaining $1/2$ & $24$\\
Pairs attaining the next ratio $1/4$ & $34$\\
\bottomrule
\end{tabular}
\end{center}
The $24$ sharp pairs are exactly the members, within the search range, of the two families
above together with $(3,4)$.  This is finite evidence only; no infinite classification is
asserted.  \computed
\end{experiment}

The sharp pairs are listed in Appendix~\ref{app:sharp-pairs}.  The finite result is useful as
a sanity check on the algebraic theorem and as evidence that the middle-monomial equality
case is highly rigid.

\begin{resultbox}
Reciprocity turns a potentially large coefficient space into a one-dimensional obstruction:
two equal-degree reciprocal filters cannot erase more than half the degree at infinity.  In
the conditional two-scale determinant, even this optimal contact leaves an exponential
integer of rate at least $D(\theta+1/2)s\log M$.
\end{resultbox}

\section{Synthesis: what the new route proves and what it rules out}
\label{sec:synthesis}

\subsection{The height/contact ledger}

The constructions above can be summarized by two quantities.  For a rational approximation
$z_j\in\Q$ to a fixed real number $\xi$, define
\begin{equation}
 \mathcal A_j(\xi)=
 \frac{-\log|z_j-\xi|}{h(z_j)}.
 \label{eq:amplification-ratio}
\end{equation}
A successful irrationality-measure contradiction requires a positive ratio exceeding the
known effective threshold for $\xi$.  The new secant families satisfy
\[
 \mathcal A_j\!\left(\log_3 2\right)\longrightarrow0.
\]
For fixed Richardson order $m$, the numerator grows from $O(\log s_j)$ to
$mO(\log s_j)$, but the denominator remains $\Theta(s_j)$.  No fixed $m$ changes the limit.

The common-residual analysis explains why denominator reduction does not secretly rescue the
construction.  The reduced height is already $\kappa s_j+o(s_j)$, and every fresh-prime gcd
is $\exp(o(s_j))$.  The only exponential cancellation occurs at the fixed structural primes.

\subsection{Families now excluded}

The following strategies cannot close the conjecture without a genuinely new ingredient.

\begin{enumerate}
\item A single normalized secant, in any fixed pair of toric directions.
\item Any fixed finite Richardson, Prouhet, or finite-difference combination of toric
secants.
\item Any attempt to obtain exponential denominator clearing from primes outside the fixed
support of $6MA$ for a fixed finite residual family.
\item Any fixed equal-degree reciprocal filter, including fixed products or quotients of
cyclotomic polynomials, evaluated along the powers $(M^s,A^s)$.
\item Any argument based on the raw numerator or common denominator before primitive
reduction.  The exact slope $\kappa$ is the relevant invariant.
\end{enumerate}

\subsection{Why this is still not a proof}

Corvaja--Zannier is a lower-height theorem outside algebraic subtori.  Our orbit avoids those
subtori, so fixed rational functions retain height.  But the theorem is not uniform over a
sequence of rational functions whose degrees, supports, and coefficients grow with $j$.
Such a sequence may have output-specific cancellations invisible to every fixed-function
statement.

Likewise, auxiliary-prime gcds are subexponential, but structural-prime valuations can be
linear and are not bounded away by the Subspace Theorem.  A proof may still exist in a
carefully designed $2$- and $3$-adic construction, a growing Hermite--Pad\'e family, or a
hybrid in which the same coefficients create both high archimedean contact and primitive
structural-prime cancellation.

\section{The surviving research target}
\label{sec:target}

\subsection{A primitive height/contact amplifier}

The secant law suggests a precise target.

\begin{target}[Growing rational observable]
\label{target:growing-observable}
Construct rational functions $F_j\in\Q(X,Y)$ and convergents $r_j/s_j\to\beta$ such that,
for $z_j=(u_j,v_j)$,
\begin{enumerate}
\item $F_j(z_j)\in\Q$ tends to a fixed logarithmic constant, preferably
$\rho=\log_3 2$;
\item the analytic contact order $k_j$ along $(2^t,3^t)$ tends to infinity;
\item after complete numerator/denominator reduction,
\[
 h(F_j(z_j))=o(k_j\log s_j),
\]
or at least is small enough that the resulting approximation beats an explicit irrationality
measure for $\rho$;
\item the height saving is supported at the structural primes and is not merely a raw common
denominator or an auxiliary-prime gcd.
\end{enumerate}
\openstatus
\end{target}

Every fixed $F$ fails condition 3 by the Corvaja--Zannier theorem.  Every fixed reciprocal
cyclotomic filter fails already at the contact stage.  Thus a surviving family must be both
growing and nontrivially primitive.

\subsection{A finite arithmetic criterion}

One can formulate a proof-producing version.  Let $F=P/Q$ with coprime integer polynomials,
and evaluate at a concrete rational point $(u,v)$ from~\eqref{eq:uv}.  Compute the exact
reduced rational $p/q=F(u,v)$ and a rigorous interval enclosure for $|p/q-\rho|$.  Any
certificate satisfying
\begin{equation}
 0<\left|\frac pq-\rho\right|<q^{-\mu_*}
 \label{eq:irrationality-certificate}
\end{equation}
with $\mu_*$ larger than a certified irrationality exponent for $\rho$ would contradict the
counterexample hypothesis.  Theorems~\ref{thm:secant-law} and~\ref{thm:richardson} show that
no fixed support can produce such certificates asymptotically.  A search should therefore
optimize reduced height directly over growing supports.

\subsection{Suggested computational search}

A practical experiment can be organized as follows.

\begin{enumerate}
\item Choose a symbolic support $\mathcal S_j\subset\Z^2$ whose size grows moderately.
\item Impose $k_j$ linear Taylor-contact equations at $(1,1)$ along
$t\mapsto(2^t,3^t)$.
\item Solve over $\Z$ or over a lattice adapted to the known $2$- and $3$-adic valuations.
\item Evaluate the resulting numerator and denominator at exact convergent points
$(u_j,v_j)$.
\item Reduce the rational number completely and record its Smith/valuation profile, not its
unreduced denominator.
\item Reject every family whose primitive height divided by $k_j\log s_j$ does not decrease.
\end{enumerate}

The report's theorems provide calibration tests: a correct implementation must reproduce the
slope $\kappa$ for the basic secant, linear height for every fixed Richardson order, and the
half-degree ceiling for reciprocal filters.

\section{Formalization blueprint}
\label{sec:lean}

The new elementary and height bookkeeping can be separated cleanly from the classical deep
inputs.

\subsection{Kernel-suitable finite core}

The following statements are suitable for direct formalization in Lean or another proof
assistant.

\begin{enumerate}
\item Definition of $u(r,s)=2^r/M^s$ and $v(r,s)=3^r/A^s$ as rational numbers; proof of
multiplicative independence when $r-s\beta\ne0$.
\item Exact $p$-adic valuation formulas and the finite-sum formula for
$h(1:u:v)$, with convergence to~\eqref{eq:kappa} under $r/s\to\beta$.
\item Taylor expansion~\eqref{eq:sigma-expansion} and the finite-difference moment identities
used in Theorem~\ref{thm:richardson}.
\item The Bernoulli coefficient formula and nonvanishing reduction to transcendence of
$\log_2 3$.
\item The reciprocal half-contact theorem, including the equality characterization.
\item The M\"obius-layer formula~\eqref{eq:cyclotomic-ratio} and the deduction
$K(n,m)\le\varphi(n)/2$.
\item A certificate checker for the finite scan: totients, M\"obius values, divisor lists,
and first unequal layer can all be verified by exact integer computation.
\end{enumerate}

\subsection{Explicit external interfaces}

Two classical results should remain named hypotheses until formal libraries are available.

\begin{description}
\item[Two-logarithm lower bound.]  For nonzero
$r\log2-s\log M$, a polynomial lower bound in $\max\{|r|,|s|\}$.
\item[Corvaja--Zannier.]  The secant-height asymptotic and the fixed rational-function lower
height outside finitely many proper torus translates.
\end{description}

All deductions from those interfaces are elementary height, valuation, and Taylor algebra.
In particular, the exact gcd slope~\eqref{eq:G-asymptotic} should be formalized as a short
corollary once the secant-height interface is assumed.

\section{Conclusion}
\label{sec:conclusion}

The new route begins with a hypothetical counterexample and follows its continued-fraction
convergents in the opposite direction from the usual determinant method.  Instead of seeking
many integer points at large scale, it compresses two adjacent multiplicative descriptions
into a rational $S$-unit point converging to $(1,1)$.  This makes the archimedean contact
excellent and the nonarchimedean height visible.

The resulting picture is exact.  The basic rational secant has primitive height
$\kappa s+o(s)$; its reduced numerator and denominator each have that rate.  The raw
cross-residual gcd has a computable complementary rate, and every part arising from new
primes is subexponential.  Fixed-order extrapolation improves only the power of a
polynomially small error.  Reciprocal cyclotomic filters lose at least half their nominal
contact before arithmetic is even considered.

This does not cross the four-exponentials frontier.  It does narrow the battlefield.  A
successful proof along the new seam must be a growing-complexity, primitive-height
construction whose cancellation is tied specifically to the structural primes and to the
conditional output pair.  Fixed rational observables, fresh-prime gcds, and reciprocal
cyclotomic filters are no longer plausible places for the missing exponential gain to hide.

\appendix

\section{Detailed height calculation}
\label{app:height}

For completeness, this appendix writes the projective height in elementary integer terms.
Let
\[
 u=\frac{U_1}{V_1},\qquad v=\frac{U_2}{V_2}
\]
be reduced positive rationals and let $D=\lcm(V_1,V_2)$.  Then
\[
 h(1:u:v)=
 \log\frac{\max\{D,U_1D/V_1,U_2D/V_2\}}
 {\gcd(D,U_1D/V_1,U_2D/V_2)}.
\]
Applied to~\eqref{eq:uv}, one may take, for large $j$,
\begin{align*}
 U_1&=2^{r_j-ms_j},&V_1&=M_0^{s_j},\\
 U_2&=3^{r_j-as_j},&V_2&=A_0^{s_j}.
\end{align*}
The three unreduced projective coordinates are all archimedean-comparable because
$u_j,v_j\to1$.  Their common content is nevertheless exponential and is exactly what the
local maximum in~\eqref{eq:kappa} removes.  This is why replacing $h(1:u:v)$ by
$\log\lcm(V_1,V_2)$ can overstate the useful height.

At each prime $p$, the primitive denominator exponent is
\[
 \max\{0,-v_p(u_j),-v_p(v_j)\},
\]
which after division by $s_j$ converges to the local summand in~\eqref{eq:kappa}.  No limiting
interchange issue occurs because only finitely many primes divide $6MA$.

\section{Bernoulli coefficient formula}
\label{app:bernoulli}

Let $a,b\ne0$ and
\[
 \frac{e^{at}-1}{e^{bt}-1}=\sum_{k\ge0}c_k(a,b)t^k.
\]
Using $z/(e^z-1)=\sum_{n\ge0}B_nz^n/n!$ gives
\begin{equation}
 c_k(a,b)=
 \frac1b
 \sum_{n=1}^{k+1}
 \frac{a^n}{n!}
 \frac{B_{k+1-n}b^{k+1-n}}{(k+1-n)!}.
 \label{eq:ck-bernoulli}
\end{equation}
The $n=k+1$ term is $a^{k+1}/(b(k+1)!)$, so the homogeneous polynomial is nonzero.  Setting
$b=\theta a$ gives
\[
 c_k(a,\theta a)=a^kR_k(\theta)
\]
for a nonzero rational function $R_k$ whose numerator is a nonzero polynomial in $\theta$.
Transcendence of $\theta$ proves $c_k\ne0$.

The finite-difference moments follow from
\[
 \sum_{i=0}^m(-1)^i\binom mi e^{it}=(1-e^t)^m.
\]
Differentiating $k$ times at $t=0$ yields zero for $k<m$ and $(-1)^mm!$ for $k=m$.

\section{Sharp cyclotomic pairs found in the scan}
\label{app:sharp-pairs}

The $24$ pairs $(n,m)$ with $3\le n<m\le50000$, equal totient, and
$K(n,m)/\varphi(n)=1/2$ were
\begin{center}
\begin{tabular}{rrrr}
\toprule
$n$ & $m$ & $\varphi(n)$ & $K$\\
\midrule
3&4&2&1\\
3&6&2&1\\
4&6&2&1\\
8&12&4&2\\
9&18&6&3\\
16&24&8&4\\
27&54&18&9\\
32&48&16&8\\
64&96&32&16\\
81&162&54&27\\
128&192&64&32\\
243&486&162&81\\
256&384&128&64\\
512&768&256&128\\
729&1458&486&243\\
1024&1536&512&256\\
2048&3072&1024&512\\
2187&4374&1458&729\\
4096&6144&2048&1024\\
6561&13122&4374&2187\\
8192&12288&4096&2048\\
16384&24576&8192&4096\\
19683&39366&13122&6561\\
32768&49152&16384&8192\\
\bottomrule
\end{tabular}
\end{center}
The ordering in this table is by increasing maximum index, not by totient.  Within the finite
range, all but $(3,4)$ are instances of
\[
 (2^a,3\cdot2^{a-1})\quad(a\ge2),
 \qquad
 (3^a,2\cdot3^a)\quad(a\ge1).
\]

\section{Finite scan source}
\label{app:code}

The following pure-Python program uses no external package.  It computes Euler totients and
M\"obius values by sieves, groups indices by totient, and compares the finite coefficient maps
in~\eqref{eq:cd}.

\begin{lstlisting}[style=py]
from collections import Counter, defaultdict
from fractions import Fraction

LIMIT = 50_000

# Euler phi sieve.
phi = list(range(LIMIT + 1))
for p in range(2, LIMIT + 1):
    if phi[p] == p:
        for n in range(p, LIMIT + 1, p):
            phi[n] -= phi[n] // p

# Linear Mobius sieve.
mu = [0] * (LIMIT + 1)
least_prime = [0] * (LIMIT + 1)
primes = []
mu[1] = 1
for n in range(2, LIMIT + 1):
    if least_prime[n] == 0:
        least_prime[n] = n
        primes.append(n)
        mu[n] = -1
    for p in primes:
        if p > least_prime[n] or p * n > LIMIT:
            break
        least_prime[p * n] = p
        mu[p * n] = 0 if n % p == 0 else -mu[n]

def divisors(n: int) -> list[int]:
    small, large = [], []
    d = 1
    while d * d <= n:
        if n % d == 0:
            small.append(d)
            if d * d != n:
                large.append(n // d)
        d += 1
    return small + large[::-1]

def layer_map(n: int) -> dict[int, int]:
    return {d: mu[n // d] for d in divisors(n) if mu[n // d] != 0}

groups = defaultdict(list)
for n in range(3, LIMIT + 1):
    groups[phi[n]].append(n)

pair_count = 0
sharp = []
ratio_counts = Counter()

for degree, indices in groups.items():
    if len(indices) < 2:
        continue
    maps = {n: layer_map(n) for n in indices}
    for i, n in enumerate(indices):
        left = maps[n]
        for m in indices[i + 1:]:
            right = maps[m]
            pair_count += 1
            keys = set(left) | set(right)
            contact = min(
                d for d in keys if left.get(d, 0) != right.get(d, 0)
            )
            ratio = Fraction(contact, degree)
            ratio_counts[ratio] += 1
            if 2 * contact == degree:
                sharp.append((n, m, degree, contact))
            if 2 * contact > degree:
                raise AssertionError((n, m, degree, contact))

print("pairs", pair_count)
print("maximum ratio", max(ratio_counts))
print("sharp count", len(sharp))
print("next ratios", sorted(ratio_counts.items(), reverse=True)[:6])
for row in sharp:
    print(*row)
\end{lstlisting}

The run used CPython 3 and completed in a few seconds on the report-generation host.  Its
headline output was
\begin{verbatim}
pairs 665095
maximum ratio 1/2
sharp count 24
next ratios [(1/2, 24), (1/4, 34), (1/6, 22),
             (1/8, 145), (1/10, 25), (1/12, 186)]
\end{verbatim}

\section{Classical inputs and bibliography}
\label{app:bib}

\begin{thebibliography}{99}

\bibitem{AlaogluErdos1944}
L.~Alaoglu and P.~Erd\H{o}s,
\newblock On highly composite and similar numbers,
\newblock \emph{Transactions of the American Mathematical Society} \textbf{56} (1944),
448--469.

\bibitem{Baker1975}
A.~Baker,
\newblock \emph{Transcendental Number Theory},
\newblock Cambridge University Press, 1975.

\bibitem{BCZ2003}
Y.~Bugeaud, P.~Corvaja, and U.~Zannier,
\newblock An upper bound for the G.C.D. of $a^n-1$ and $b^n-1$,
\newblock \emph{Mathematische Zeitschrift} \textbf{243} (2003), 79--84,
\newblock DOI: \url{https://doi.org/10.1007/s00209-002-0449-z}.

\bibitem{CZ2005}
P.~Corvaja and U.~Zannier,
\newblock A lower bound for the height of a rational function at $S$-unit points,
\newblock \emph{Monatshefte f\"ur Mathematik} \textbf{144} (2005), 203--224,
\newblock arXiv: \url{https://arxiv.org/abs/math/0311030},
DOI: \url{https://doi.org/10.1007/s00605-004-0273-0}.

\bibitem{LMN1995}
M.~Laurent, M.~Mignotte, and Yu.~Nesterenko,
\newblock Formes lin\'eaires en deux logarithmes et d\'eterminants d'interpolation,
\newblock \emph{Journal of Number Theory} \textbf{55} (1995), 285--321.

\bibitem{Matveev2000}
E.~M.~Matveev,
\newblock An explicit lower bound for a homogeneous rational linear form in logarithms of
algebraic numbers. II,
\newblock \emph{Izvestiya: Mathematics} \textbf{64} (2000), 1217--1269.

\bibitem{Waldschmidt2000}
M.~Waldschmidt,
\newblock \emph{Diophantine Approximation on Linear Algebraic Groups},
\newblock Springer, 2000.

\bibitem{Yasufuku2026}
Y.~Yasufuku,
\newblock GCD inequalities arising from codimension-2 blowups,
\newblock \emph{Bulletin of the London Mathematical Society} \textbf{58} (2026), e70349,
\newblock DOI: \url{https://doi.org/10.1112/blms.70349}.

\bibitem{AnthropicFable2026}
Anthropic,
\newblock Claude Fable 5 and Claude Mythos 5,
\newblock 9 June 2026,
\newblock \url{https://www.anthropic.com/news/claude-fable-5-mythos-5}.

\bibitem{AnthropicCrypto2026}
Anthropic,
\newblock Discovering cryptographic weaknesses with Claude,
\newblock 28 July 2026,
\newblock \url{https://www.anthropic.com/research/discovering-cryptographic-weaknesses}.

\end{thebibliography}

\end{document}
